% CEUR-WS single-column manuscript, per https://ceur-ws.org/HOWTOSUBMIT.html#CEURART
% Requires ceurart.cls (download the CEURART bundle and place ceurart.cls in the
% project root, e.g. from the Overleaf "CEURART" template).
% Compile with pdfLaTeX + BibTeX. Prose avoids dash punctuation by request.
% Statistics use math superscripts. Units are written as text ($\mu$M, nM).
\documentclass[sigconf]{ceurart}

\usepackage{amsmath}
\usepackage{booktabs}

\begin{document}

\copyrightyear{2026}
\copyrightclause{Copyright for this paper by its authors.
  Use permitted under Creative Commons License Attribution 4.0
  International (CC BY 4.0).}

\conference{18th International SWAT4HCLS Conference, 2026. Basel, Switzerland (SWAT4HCLS)}

\title{When explicitness pays in a tool's return value:
a controlled study on a production life-science MCP server}

\author[1]{Yasunori Yamamoto}[email=yy@dbcls.rois.ac.jp,orcid=0000-0002-6943-6887]
\author[1]{Yuki Moriya}[email=moriya@dbcls.rois.ac.jp,orcid=0000-0001-8195-5893]
\address[1]{Database Division for Life Science (DBCLS), National Institute of Genetics, 178-4-4 Wakashiba, Kashiwa, Chiba 277-0871, JAPAN}
\author[2]{Akira R. Kinjo}[email=arkinjo@gmail.com,orcid=0000-0002-4006-8208]
\address[2]{Anima Machina G.K., Osaka Station Building No. 3, 29th floor, Room 1-1-1, 1-1-3 Umeda, Kita-ku, Osaka 530-0001, Japan}

\begin{abstract}
Life-science databases are increasingly reached by AI agents through a shared
interface called the Model Context Protocol (MCP). Existing advice tells whoever
builds one to always be explicit: state the units, spell out identifiers, label
each column. It never says where the added clarity is worth the tokens it costs,
or for which agent. We studied this on a production server, holding the model and
task fixed and changing only one column of the returned payload, then measuring
how the answer changed and at what token cost.

Whether being explicit pays depended on which agent was asking, not on what the
data contained. In the first setting the information is present but implicit: a
compound's label is either its real name or a copy of its own identifier code,
which also appears in the next column. A higher-performing model resolves this
itself (123 of 124 correct), so a column stating the answer helps it almost
nothing. A lower-performing model mistakes the code for a name; adding one column
marking each label as name or code, data it already received, raises it from 16 to
75 of 124 for about 407 tokens, under one percent of the run.
Even so, the latter still falls short of the former's performance.

In the second setting the information is truly missing: with the unit removed, a
bare 10 could mean nanomolar or micromolar, and no reasoning recovers which. Across
twenty targets this made the higher-performing model wrong on 350 of 929
micromolar rows and the lower-performing model wrong on all, while correct rows
still cost several times more. Explicitness is best understood as a safeguard: a few
hundred tokens that do nothing for an agent that would infer the answer and a great
deal for one that would not, applied to every client because a server cannot tell in
advance which one needs it.
\end{abstract}

\begin{keywords}
  Model Context Protocol \sep
  LLM agents \sep
  interface design \sep
  return-value explicitness \sep
  token cost
\end{keywords}

\maketitle

\section{Introduction}

Connecting a language model to knowledge it does not hold internally has been a
standing requirement rather than a passing one, whether the connection is made by
retrieval~\cite{lewis2020rag}, by learned tool use~\cite{schick2023toolformer}, or
by the broader family of methods that augment a model with external
resources~\cite{mialon2023augmented}. AI agents that use language models
increasingly reach databases not through custom code but through a thin, shared
protocol. In the life sciences, the Model Context Protocol (MCP)~\cite{mcp2024spec}
is quickly becoming that layer. By early 2026, it had more than 10{,}000 active servers
and tens of millions of monthly SDK downloads~\cite{srinivasan2026bridging}
across all domains. Even within the life sciences alone, a curated list we maintain
records several dozen hosted remote MCP servers from independent providers, spanning
public institutions such as EMBL-EBI, SIB, and DBCLS as well as commercial and
community implementations~\cite{dbcls2026remotemcp}. A server advertises a few
tools, an agent calls them,
and whatever the tool returns becomes the agent's view of the data.
The person who builds the server therefore decides, in the shape of each response,
how much of the reasoning is done for the agent and how much is left to it.

Guidance for that builder is plentiful and consistent. Limit the size of
responses, paginate long results, return errors as structured values rather than
prose, and above all be explicit: name the units, write out the identifiers, and
say what each column means. Reports from production MCP deployments turn these same
concerns, such as server contracts and structured error semantics, into readiness
checklists~\cite{srinivasan2026bridging}. The advice is sound as far as it goes,
but it is uniform. It says to be explicit everywhere, and it does not say where
that explicitness earns back the tokens it costs, or for whom. On a public
endpoint, those tokens are paid on every call, by every client, indefinitely.

We treat that gap as an empirical question. Keeping the model and the task fixed
and changing only the payload that a production server returns, using
TogoMCP~\cite{kinjo2026togomcp} over the EBI RDF platform's~\cite{jupp2014ebirdf}
ChEMBL graph~\cite{zdrazil2024chembl}, we ask when added explicitness changes an
agent's answer, and at what cost. We keep the less explicit version realistic. In
the label experiment it is not invented at all: it is the data as the server already
stores it, because a compound's \texttt{rdfs:label} is its accession rather than a
real name for 97.7 percent of ChEMBL molecules (Section~\ref{sec:results}). In the
units experiment we instead reconstruct a poorer design by removing a column that ChEMBL
itself provides. ChEMBL does record the unit, so the omission is not a fault of
ChEMBL; it reproduces the case of a simpler interface or export that returns a bare
value with no unit. This is not an exotic case: any query or export that selects the
value column without the unit column yields exactly this payload, including an agent
that composes its own SELECT and omits \texttt{?units}. The two experiments play
complementary roles. The label experiment is grounded in real data at a measured prevalence, because it
uses the server's own data as stored. The units experiment is a controlled
manipulation that isolates the case of genuinely absent information, which no
amount of reasoning can recover; we read it as a clean demonstration of that
mechanism, not as a claim about how often this particular column is dropped in
practice. Two settings organize the study. Both begin from a query that
asks the server for the compounds tested against a chosen protein, which in this
field is called a target, and returns one row per compound. In the first setting
everything the agent needs is present but implicit: a compound's
\texttt{rdfs:label} is either a real name or a copy of the compound's own
accession, and that accession sits in the next column, so the ambiguity can be
resolved from the payload alone. In the second setting the needed information is
genuinely absent: with the unit column removed, a bare 10 could mean 10 nanomolar
or 10 micromolar, and no reasoning over the returned rows recovers which.

Our central finding is that whether being explicit helps depends on which agent is
reading the response, not on what the response contains. For a lower-performing
agent the explicit column helps in both settings, but its role differs. When the
information is only implicit, the agent still succeeds sometimes without it (16 of
124), and the column raises its score sharply without letting it catch up. When
the information is genuinely absent, the column is necessary, because without the
unit the agent succeeds on none. The higher-performing agent, by contrast,
resolves the implicit case on its own (123 of 124), so the column buys it almost
nothing there, while the absence of a genuinely needed unit still damages it and
raises its cost. A server cannot see which agent is calling, so it cannot apply
this cost only to the clients that need it, which is why a uniform instruction to
be explicit is defensible as a precaution and hard to defend as efficiency advice.

The effects rest on patterns that are common rather than anecdotal. A compound's
label is its accession for 97.7 percent of ChEMBL's 1.9 million small molecules,
for 98.4 percent of the molecules that bind TP53, and, in a second database, for
9.8 percent of HomoloGene's genes. One point about the method is worth stating
here, because a broken change in this setting does not crash but instead returns
fluent, plausible prose. We therefore check every number by reading the actual
bytes that the server returned and the agent wrote, rather than by trusting a
summary. We mentioned the specific cases in which this
discipline reversed our own early conclusions in Discussion.

We make four contributions. First, to our knowledge the first controlled
measurement of the explicitness of a server's return value on a production MCP
server, holding the model and the task fixed, together with a reusable measurement
framework and task suite. Second, the finding that the value of that change
depends on the client: for a lower-performing agent the explicit column is highly
beneficial but not equalizing when the information is only implicit, and necessary
when the information is genuinely absent, while for a higher-performing agent it is
unneeded in the first case and only a cost in the second. Third, a cost account
with clear denominators and an asymmetry in which a withheld column raises token
cost on rows it never corrupts, the same amplification that earlier work produces
on purpose through attacks, here appearing from an ordinary, less explicit payload.
Fourth, evidence that the underlying patterns are the common case, and a method for
measuring interface effects when a broken change returns fluent but wrong prose.

\section{Related Work}

\paragraph{Benchmarking tool use under MCP.}
One line of work measures how well models use tools. MCP-RADAR scores 507 tasks
across six domains on answer correctness and operational accuracy, and reports a
trade-off between accuracy and efficiency across models~\cite{gao2025mcpradar}.
MCP-Bench and MCPToolBench++ extend this to hundreds and thousands of live
servers~\cite{wang2025mcpbenchbenchmarkingtoolusingllm,fan2025mcptoolbenchlargescaleai}. These studies vary the model
and the task over fixed tools in order to profile capability. We invert that
design. We hold the model and the task fixed and vary one column of the server's
output, so that any effect can be attributed to that design choice rather than to a
model's aptitude, and we ground the study in a well built production server rather
than a benchmark harness.

\paragraph{Designing MCP servers and tools.}
A second line prescribes how to build the server. Field lessons from an enterprise
MCP deployment group failures into five design areas, namely server contracts,
user context, timeouts, errors, and observability, and propose structured error
semantics and adaptive tool budgeting~\cite{srinivasan2026bridging}. On the tool
description side, an audit of 856 tools finds that 97.1 percent carry at least one
defect and shows that enriching descriptions raises success by a median of 5.85
points but adds 67 percent more execution steps and makes 17 percent of cases
worse~\cite{hasan2026smelly}. We share the message that explicitness is not free
and can even hurt, but we move it from the input side, meaning how a tool is
described, to the output side, meaning what a tool returns. We add the client
dependence that a single agent study cannot see, and the column we add can be
computed from two columns the response already contains, so it gives the agent no
new information. Any improvement therefore comes from stating what was already
present, not from providing something new.

\paragraph{The cost of tool interaction.}
A third line treats the token cost of tool use as its object of study and trims
verbose tool output to reduce it~\cite{xiao2026reducingcostllmagents}. Most directly, an
economic denial of service attack acts only on the content an MCP tool returns to
the model, the text the agent reads, while leaving the tool's declared interface
(its function signatures) and the final answer it produces unchanged. By
restructuring that returned content, it steers the agent into long tool calling
chains that raise the cost per query by as much as 658 times, yet the task still
completes correctly, which makes the tampering hard to
detect~\cite{zhou2026beyondmaxtokens}. Our units experiment acts on the same surface, the
returned content rather than the declared interface, but from the opposite
direction: an ordinary, benign, less explicit payload, not an engineered attack. It
reproduces the same correctness independent cost inflation and, unlike the attack,
also causes a selective loss of correctness on exactly the rows the omission
corrupts. The cost increase that an attacker engineers, a careless design incurs for
free.

\paragraph{When to call a tool.}
A fourth line studies the decision of whether to call a tool at all. When2Tool
shows that the need for a tool call can be read from a model's hidden state (AUROC
0.89 to 0.96), that models already know when a call is needed but fail to act on
that knowledge, and that prompting cannot cleanly control the balance between
accuracy and the number of calls~\cite{sun2026llmagentsknowtools}. Their concern is calling
too often and reducing unnecessary calls. Ours is the opposite. By withholding a
column, the server creates the need for an extra call, and in our data the agents
that recovered were largely those that noticed the gap and paid for that extra call
(Section~\ref{sec:results}). Their finding that prompting is a blunt lever supports
our claim that the effective lever sits in the server's response, not in the
client's prompt.

\paragraph{Context engineering and model capacity.}
A fifth line is closest to our result about the client. It shows that carefully
built tool environments let small models match large ones. MARRVEL-MCP provides
models of 3 to 20 billion parameters with a set of 44 curated genomics tools and
reports that a 20 billion parameter model rises from 41 to 94 percent, concluding
that contextual guidance can compensate for limited model
capacity~\cite{everton2026marrvelmcp}. The same spirit appears in agent frameworks
that curate or evolve their tool sets. Their result is itself quantitative, but it
measures the effect of a whole curated tool environment. It does not separate the
value of a single explicitness choice, report a case where such a choice is worth
nothing, or show that the help can be large yet still leave a lower-performing
model behind a higher-performing one. We add these three things, and conclude that
contextual guidance compensates for limited capacity only in part, not in full.

\section{Materials and Methods}

\paragraph{Server, data, and tool.}
All queries were issued to a production MCP server, TogoMCP~\cite{kinjo2026togomcp},
backed by the EMBL EBI RDF platform~\cite{jupp2014ebirdf}, over the ChEMBL (version
34) graph~\cite{zdrazil2024chembl}. The server exposes a single
SPARQL tool, \texttt{run\_sparql}, which an agent calls with a database name and a
query string and which returns rows as CSV. We never modified the server's query
engine. We varied only the shape of the CSV it returned, using a reshaping layer
placed in front of the tool that, per condition, either passed the native result
through or removed or added one column. Every task instructs the agent to run a
single fixed (pinned) query verbatim, so any difference between conditions arises
from reshaping rather than from a query the agent composed. We also use a second
resource, HomoloGene~\cite{ncbi2013resources}, but only to count how often its label
column carries a bare identifier rather than a name; no agent was run against
HomoloGene.

\paragraph{Models and decoding.}
Two open weight agents were evaluated: Qwen3.6-35B-A3B (the higher-performing
agent) and Qwen3-30B-A3B (the lower-performing agent), both from the Qwen3
family~\cite{qwen3technical}. Both were served locally
with vLLM~\cite{kwon2023vllm} at temperature 0, with a 16{,}384 token generation
reserve, thinking disabled, a recursion limit of 60, and a single sequential
worker. The two models require different serving profiles; in particular, the
lower-performing model needs a YaRN~\cite{peng2023yarn} rope scaling factor of 2.0
and a maximum context of 65{,}536, and a wrong
factor silently corrupts generation, so the serving command line was verified
before every run. In our runs both models were effectively deterministic at
temperature 0: repeated runs gave identical outputs. Minor variation from software
or hardware is possible in principle, but we did not observe it, which is why one
run per item per condition is sufficient and why any variation across items
reflects the items rather than sampling. Throughout, an item is one task: a single
question about one molecule in the label experiment or one compound in the units experiment, run
once under each condition. When we report tokens for a run, we mean the total
number of tokens the agent used for one item under one condition, summed over the
prompt and the generation across every step and tool call of that run.

\paragraph{Two manipulations.}
We study two settings, each a single column change to the returned payload. In the
\emph{label experiment} (recoverable ambiguity) the information the agent needs is present
but implicit. For a chosen target we ask the server for the compounds tested
against it, each with its \texttt{rdfs:label} and \texttt{chemblId}. A molecule's
label is either a real compound name or a copy of the molecule's own accession, and
the accession also appears in the adjacent column, so the ambiguity is resolvable
from the payload alone. The two conditions are \texttt{label\_implicit}, the native
table, and \texttt{label\_explicit}, the same table plus one column,
\texttt{moleculeLabel\_kind}, whose value is \texttt{id} when the label equals the
accession and \texttt{name} otherwise. That column is a deterministic function of
two columns present in both conditions, so it adds no information and only makes the
distinction explicit. Each task names one molecule and asks the agent to report its
real compound name, or exactly NONE if it has none. In the \emph{units experiment}
(genuinely absent information) we ask for the IC50 activities of a target's
compounds as \texttt{chemblId}, \texttt{value}, and \texttt{units}, filtered to
rows whose unit is nanomolar (nM) or micromolar ($\mu$M). The two conditions are
\texttt{units\_on}, the native table, and \texttt{units\_off}, the same table with
the units column removed. With the unit removed, the default reading of a bare
value is nM, which is correct for a row whose true unit was nM and wrong by a factor
of 1000 for a row whose true unit was $\mu$M. Each task names one compound and asks for
its IC50 in nM.

\paragraph{Target selection.}
The label experiment uses a single, canonical target, CHEMBL4096 (the tumor protein TP53),
because the manipulation there tests a data curation pattern that is a property of
the identifier, not of the compound, and because generality is carried by the
exhaustive prevalence counts. The units experiment uses a random sample of targets,
because there the effect could in principle be a property of one target's
particular values. The sampling population is the set of ChEMBL targets that have
both nM and $\mu$M units recorded for IC50, which is a scope condition (a manipulation
of the unit column is only defined where both units occur), not a selection on the
size of any effect. That population contains 3{,}629 targets. From it we drew 20
targets by a seeded, order independent rule, $\mathrm{rank}(id) =
\mathrm{sha256}(seed, id)$ ascending, so the sample is reproducible and does not
depend on the query engine's row order.

\paragraph{Deterministic response windows.}
Each pinned query returns at most a fixed number of rows (LIMIT 200 for the label
experiment, LIMIT 100 for the units experiment). A LIMIT without an explicit ordering returns an
implementation defined subset, which is not reproducible when the number of
candidate rows exceeds the limit; we observed this directly (repeated executions
returned different subsets). We therefore added a deterministic ORDER BY to every
pinned query: ORDER BY \texttt{?chemblId} for the label experiment and ORDER BY
\texttt{?activity} for the units experiment. We verified that \texttt{?activity} is unique
per returned row for all 20 units targets (the row count equals the count of
distinct activities), so it induces a total order and a reproducible window;
\texttt{?chemblId} plays the same role for the label experiment. We build each
target's item list, with its frozen correct and trap answers, from our own run of
the pinned query; at test time the agent runs the same query again through
\texttt{run\_sparql}. Scoring is therefore valid only if the two runs return the
same rows, which is why the ordering must be deterministic. With the ordering fixed,
the item selection window, the agent's live query, and any later capture of the
payload all return the identical rows.

\paragraph{Item construction.}
On the deterministic TP53 window of 200 rows (124 id-only and 76 named), the label
experiment uses every row: 124 treatment items (label equals accession, correct answer
NONE) and 76 control items (label is a real name, correct answer the name). For the
units experiment, for each target we applied a frozen exclusion set to the window before
selecting scored items. Two flags exclude a row: \texttt{dataValidityIssue = true}
(a pathological value that would mislead the agent before any unit reasoning) and
\texttt{potentialDuplicate = true} (a non independent repeat of another row's
value). The \texttt{standardRelation} flag (for example ``\texttt{>}'') and
\texttt{activityComment} were examined and not excluded, because they are orthogonal
to unit handling. The agent still sees the full raw window; the exclusion set only
governs which rows are scored. Scored treatment items are the eligible $\mu$M rows and
scored control items are the eligible nM rows, taken in payload order with no
selection on outcome. No per target cap was applied, so a target contributes as
many items as its window affords; the analysis accounts for this. The resulting
suite has 929 treatment ($\mu$M) items and 308 control (nM) items across the 20 targets.

\paragraph{Scoring.}
In the label experiment, correctness is judged from the final answer text by a regular
expression: a treatment item is correct if the answer states NONE; a control item
if it states the molecule's name. A treatment item is a misread if the answer
asserts the accession as the name. An answer that explicitly declines is scored as
an abstention, and a correct answer is never counted as an abstention or a misread.
In the units experiment, because the sampled $\mu$M values are mostly decimals and some are in
exponential notation, digit aligned string matching does not generalize, so scoring
is numeric. From the answer text we take the numbers the agent commits to and
normalize each to nM, recognizing both Unicode code points that render as the micro
sign (U+00B5 MICRO SIGN and U+03BC GREEK SMALL LETTER MU), which agents use
interchangeably. We then ask which hypothesis each number supports: the correct
value (the $\mu$M value times 1000) or the trap value (the $\mu$M value read as nM).
The two are always a factor of 1000 apart. A single unit bearing number is classified by which
value it is nearer in $\log_{10}$ space, with the boundary at their geometric mean;
a value farther than 0.5 in $\log_{10}$ from the nearer hypothesis is scored as
neither. When an answer mentions both hypotheses, the committed value is taken from
the last sentence that states a concentration, and an answer whose final sentence
explicitly declines to choose is routed to human review rather than guessed. The
gate for correctness over misread over abstention is unchanged. For control items
there is no trap. Every scoring rule was checked against the raw answer text, which
exposed two extractor defects. One micro sign code point was not recognized. The
other rule sometimes picked up a number the answer mentioned only in passing, such
as another row's value, instead of the value the agent gave as its final answer.
Both were fixed, and the small number of answers routed to review were adjudicated
by hand.

\paragraph{Prevalence.}
To establish that the collisions are common rather than anecdotal, we counted them
exhaustively rather than estimating from a sample: the fraction of ChEMBL small
molecules whose \texttt{rdfs:label} equals the accession, the same fraction among
molecules tested against TP53, and the fraction of HomoloGene genes whose label is
an LOC identifier.

\paragraph{Verification and gates.}
A broken manipulation in this setting returns fluent, plausible prose rather than
an error, so we did not trust aggregates. Before accepting a target's numbers we
applied three checks per item. First, the first \texttt{run\_sparql} output had to
differ between the two conditions; identical hashes would mean the manipulation
never reached the payload. Second, the rows the agent received under
\texttt{units\_on} had to match the rows we had selected that target's items from,
in the same order; a nondeterministic window would otherwise make the two disagree.
Third, the set of scored items had to match the frozen item list. Endpoint faults that produced empty or failed queries were detected by
these checks, and the affected targets were re-run.

\paragraph{Statistics.}
The primary test in each experiment is a stratified McNemar exact test on the paired
conditions, pooled across targets by the Cochran Mantel Haenszel procedure, two
sided. The effective sample is the set of targets (paired within a target,
independent across targets); we do not treat the items as independent. Because no
per target cap was applied, larger targets contribute more items, so we also report
the per target discordant counts and note that the direction is unanimous across
targets. Specificity, a comparison between different molecules ($\mu$M failure versus
nM failure), uses Fisher's exact test. Prevalence figures are exhaustive counts,
not tests.

\section{Results}
\label{sec:results}

\paragraph{Prevalence.}
A molecule's \texttt{rdfs:label} is a copy of its own accession, rather than a real
name, for $1{,}876{,}910$ of $1{,}920{,}603$ ChEMBL small molecules (97.7 percent).
Among the $36{,}494$ molecules tested against TP53 that carry a label, $35{,}910$
(98.4 percent) have a label that is only the accession rather than a real name. In
a second, unrelated resource, HomoloGene, the same label column carries two kinds
of content: a real gene symbol for most genes, and a bare LOC identifier plus the
gene's own number when no symbol is assigned, which happens for $27{,}066$ of
$275{,}237$ genes (9.8 percent). We use HomoloGene only as an independent instance
of one column standing in for both a name and an identifier; we did not run agents
on it.

\paragraph{Label experiment: recoverable ambiguity.}
On the deterministic TP53 window (124 treatment, 76 control), the higher-performing
agent resolves the collision essentially on its own, and the annotation buys it
almost nothing; the lower-performing agent fails badly without the annotation and
improves sharply with it, but does not catch up (Table~\ref{tab:label}).

\begin{table}[t]
\centering
\caption{Label experiment on the deterministic TP53 window.}
\label{tab:label}
\begin{tabular}{llcc}
\toprule
model & condition & treatment correct & control correct \\
\midrule
Qwen3.6-35B-A3B & implicit & 123/124 & 76/76 \\
Qwen3.6-35B-A3B & explicit & 124/124 & 76/76 \\
Qwen3-30B-A3B   & implicit & 16/124  & 76/76 \\
Qwen3-30B-A3B   & explicit & 75/124  & 76/76 \\
\bottomrule
\end{tabular}
\end{table}

For the lower-performing agent the annotation raises treatment correctness from
16/124 to 75/124 (McNemar exact, 61 items gained, 2 lost, two sided $p \approx
4.4\times10^{-16}$). It does not equalize the two agents: the between model gap,
measured as items the higher-performing agent answers and the lower-performing does
not, narrows from 107 under the implicit payload ($p \approx 1.2\times10^{-32}$) to
49 under the explicit payload ($p \approx 3.6\times10^{-15}$), but remains large and
highly significant. Controls are answered correctly by both models in both
conditions, so the effect is specific to the ambiguous treatment rows. The
annotation adds 407 tokens per run. The implicit run averaged $43{,}189$ tokens for
the higher-performing agent and $129{,}146$ for the lower-performing agent, so 407
tokens is 0.94 percent and 0.31 percent of the run respectively.

\paragraph{Units experiment: genuinely absent information.}
Across the 20 targets the suite has 929 treatment ($\mu$M) and 308 control (nM) items.
Removing the unit column leaves the higher-performing agent's answers almost
unchanged on control rows and collapses them on treatment rows
(Table~\ref{tab:units}).

\begin{table}[t]
\centering
\caption{Units experiment across 20 targets.}
\label{tab:units}
\begin{tabular}{llcc}
\toprule
model & condition & treatment ($\mu$M) correct & control (nM) correct \\
\midrule
Qwen3.6-35B-A3B & units\_on  & 929/929 & 308/308 \\
Qwen3.6-35B-A3B & units\_off & 579/929 & 307/308 \\
Qwen3-30B-A3B   & units\_on  & 260/929 & 305/308 \\
Qwen3-30B-A3B   & units\_off & 0/929   & 306/308 \\
\bottomrule
\end{tabular}
\end{table}

For the higher-performing agent, removing the unit costs 350 of 929 treatment
items, with no item recovered in the other direction (stratified McNemar/CMH, 350
lost, 0 gained, two sided $p < 10^{-77}$), and the loss appears in every one of the
20 targets (Table~\ref{tab:pertarget}). The directional prediction holds: under
\texttt{units\_off} the higher-performing agent fails on 350/929 (37.7 percent) of
$\mu$M rows but only 1/308 (0.3 percent) of nM rows (Fisher exact, odds ratio
$\approx 186$, $p \approx 5\times10^{-36}$).

The cost is not selective. For the higher-performing agent, a $\mu$M item averaged
$35{,}273$ tokens under \texttt{units\_on} and $197{,}819$ under \texttt{units\_off},
a factor of 5.6; a nM item averaged $34{,}983$ and $143{,}120$, a factor of 4.1. The
nM rows are answered correctly with or without the unit yet still cost 4.1 times
more without it. Whether a corrupted row is recovered is strongly associated with
whether the agent noticed the gap and searched. The number of calls was chosen by the agent, not assigned by us. Because we only
grouped rows by this self-selected behavior, the two groups may differ in other ways
as well, so the link between noticing and recovery is correlational rather than
causal. The agent made a second call on 766 rows and
answered in one call on 163, and was correct on 566 of the 766 but on only 13 of
the 163 (Fisher exact, $p \approx 4\times10^{-58}$). The second-call rows averaged
$232{,}473$ tokens against $34{,}967$ for single-call rows, about 6.6 times more
(Table~\ref{tab:mediation}).

The lower-performing agent shows the same shape from below. Shown the unit, it
converts $\mu$M to nM on 260/929 rows (28 percent); with the unit removed it succeeds
on none (McNemar, 260 lost, 0 gained, $p \approx 2\times10^{-58}$). The between
model gap under \texttt{units\_on} is 669 $\mu$M items that the higher-performing agent
answers and the lower-performing does not, with none in the other direction. A
single-target pilot had suggested the lower-performing agent could not convert at
all. At scale it converts on more than a quarter of rows, so the pilot's zero was a
property of that one table, not of the model.

\begin{table}[t]
\centering
\small
\caption{Units experiment, per target, higher-performing agent (Qwen3.6-35B-A3B).
Treatment items are $\mu$M rows; columns \emph{on} and \emph{off} give the number
correct under \texttt{units\_on} and \texttt{units\_off}. $b$ counts items correct
with the unit and wrong without it, $c$ the reverse. \emph{nM items} is the number
of control (nM) items for that target, and \emph{nM off} is how many of them are
correct under \texttt{units\_off}. The effect is present in every target ($b>0$)
and never reverses ($c=0$).}
\label{tab:pertarget}
\begin{tabular}{lrrrrrrr}
\toprule
target & $\mu$M items & on & off & $b$ & $c$ & nM items & nM off \\
\midrule
CHEMBL612851 & 100 & 100 & 80 & 20 & 0 & 0 & 0 \\
CHEMBL2889 & 93 & 93 & 59 & 34 & 0 & 0 & 0 \\
CHEMBL3112376 & 88 & 88 & 42 & 46 & 0 & 1 & 1 \\
CHEMBL3108638 & 84 & 84 & 63 & 21 & 0 & 1 & 1 \\
CHEMBL612570 & 73 & 73 & 38 & 35 & 0 & 0 & 0 \\
CHEMBL1075533 & 71 & 71 & 48 & 23 & 0 & 0 & 0 \\
CHEMBL614517 & 67 & 67 & 40 & 27 & 0 & 0 & 0 \\
CHEMBL1075538 & 65 & 65 & 38 & 27 & 0 & 0 & 0 \\
CHEMBL4523988 & 46 & 46 & 29 & 17 & 0 & 13 & 13 \\
CHEMBL1875 & 43 & 43 & 25 & 18 & 0 & 43 & 43 \\
CHEMBL2366254 & 43 & 43 & 28 & 15 & 0 & 11 & 11 \\
CHEMBL1743121 & 34 & 34 & 19 & 15 & 0 & 2 & 2 \\
CHEMBL2095172 & 30 & 30 & 17 & 13 & 0 & 14 & 14 \\
CHEMBL6080 & 27 & 27 & 15 & 12 & 0 & 37 & 37 \\
CHEMBL3072 & 26 & 26 & 24 & 2 & 0 & 74 & 73 \\
CHEMBL5084 & 12 & 12 & 4 & 8 & 0 & 2 & 2 \\
CHEMBL5198 & 11 & 11 & 5 & 6 & 0 & 5 & 5 \\
CHEMBL5705 & 10 & 10 & 2 & 8 & 0 & 78 & 78 \\
CHEMBL3071 & 3 & 3 & 2 & 1 & 0 & 2 & 2 \\
CHEMBL3799 & 3 & 3 & 1 & 2 & 0 & 25 & 25 \\
\midrule
TOTAL & 929 & 929 & 579 & 350 & 0 & 308 & 307 \\
\bottomrule
\end{tabular}
\end{table}

\begin{table}[t]
\centering
\small
\caption{Noticing and recovery (units experiment, treatment rows, \texttt{units\_off},
higher-performing agent). Rows are split by the agent's own number of calls. Cost
is the mean total tokens divided by the \texttt{units\_on} baseline (35{,}273
tokens).}
\label{tab:mediation}
\begin{tabular}{lrrrrrrr}
\toprule
group & $n$ & correct & misread & abstain & neither & mean tokens & cost \\
\midrule
made a second call & 766 & 566 (73.9\%) & 142 (18.5\%) & 56 (7.3\%) & 2 & 232{,}473 & 6.6$\times$ \\
answered in one call & 163 & 13 (8.0\%) & 140 (85.9\%) & 10 (6.1\%) & 0 & 34{,}967 & 1.0$\times$ \\
\bottomrule
\end{tabular}
\end{table}

\paragraph{Summary of the two experiments.}
The two experiments differ in how much the lower-performing agent needs the explicit
column. When the information is present but implicit, a zero information annotation
lifts it from 16 of 124 to 75 of 124, a large gain that still leaves it behind the
higher-performing agent. When the information is genuinely absent, the unit is
necessary for any success at all, since it succeeds on none without the unit and on
28 percent with it. In both settings the explicit column helps the lower-performing
agent far more than the higher-performing one, which needs it not at all in the
first setting and pays for its absence in the second.

\section{Discussion and Future Work}

\paragraph{Why the contribution outlives the two models we could run.}
Available hardware limited us to two open weight agents, so our capability axis has
only two points. It is worth saying plainly why this does not date the work. Our
claim is not a pair of scores. It is a structure that model progress does not
remove. A public endpoint serves many kinds of client at once: older models, small
or on-device models, quantized or low-cost deployments, and agents from other
vendors. It cannot see which one is calling. As long as that
range has any spread in capability, the value of a change to the payload depends on
the client, and the server cannot apply this cost client by client. Improving models
do not remove this spread. They lengthen the tail of lower-performing
clients that keep reaching the same endpoints for years. The safeguard logic is
therefore structural, not a feature of these particular models.

Two of our findings are, at their source, independent of the model. Information
that was never recorded, such as a missing unit, is absent from the payload itself.
No agent can recover it, no matter how capable, because capability cannot
reconstruct data that was never there. The second is that the collisions that drive
the effects, such as an accession standing in for a name or an LOC gene identifier,
are properties of the data and its curation rather than of the model. They are common, appearing for 97.7 percent of ChEMBL molecules,
98.4 percent of the molecules that bind TP53, and 9.8 percent of HomoloGene genes,
and they will remain common as models change.

\paragraph{The question outlives the protocol.}
We ran this study on an MCP server because MCP is the layer that life-science
agents use today, but the question we ask is not about MCP. How MCP itself will
develop, and whether some later standard will replace it, is not something we can
forecast. What is safe to assume is that agents built on language models will keep
needing to reach knowledge they do not hold internally, whether through retrieval,
through learned tool use, or through whatever interface succeeds MCP. The object of
our measurement is the efficiency of knowledge sharing between a language model and
an external source, meaning how much of the interpretation the source spells out
and how much it leaves to the agent, and at what cost in tokens. That trade-off,
and our central result that its value depends on the client rather than on the
payload, is a property of the exchange, not of the protocol that carries it. Any
successor interface that returns structured results to a range of agents inherits
the same choice and the same client dependence. We therefore expect the framework,
and the questions it makes measurable, to transfer to whatever replaces the
particular protocol studied here.

\paragraph{What the community can reuse.}
The lasting artifact is the measurement, not its current answer. The task sets are
small, with 200 label items and about 1{,}200 unit items. Their response windows
are fixed by ordering, and they are scored by checked criteria, so any new model
can be run through them in minutes. This turns the question of whether explicitness
is worth it for a given model into a quick, reproducible check, and it lets the
community follow a moving threshold: the capability at which a column that carries
no information stops changing the answers. If future models, smaller and more
capable, come to resolve these collisions on their own even in low-cost deployments,
the value of the added column falls to zero everywhere. The efficiency argument then
wins, which is simply not to spend tokens stating what every client already infers.
Our framework is exactly what detects that crossover. It tells a server operator
when it is safe to stop spending those tokens, and that decision support becomes more
useful, not less, as models improve.

More broadly, the method transfers beyond this server and these columns. It makes a
single change to a production payload and grounds the effect in exhaustive counts
rather than anecdote. It verifies every number by reading the bytes the server
returned and the agent wrote, because a broken change here returns fluent but wrong
prose. It fixes the response window so runs are reproducible, reports cost with
clear denominators, and reassigns failures from the model to the interface. These
pieces are reusable for any study of response design.

\paragraph{Limitations and open directions.}
The clearest limitation is the capability axis, which has only two points. Both
agents are from the Qwen family and differ in both size and generation, so calling
one higher-performing and the other lower-performing confounds capability with
generation and with family-specific behavior. We therefore report two specific
agents and do not claim a general law relating capability to the value of
explicitness. The released task sets are meant to let others place many more
agents, from other families and generations, on this axis, and to draw the full
curve that relates capability to the value of a given change. Other directions
include other servers and domains, other payload factors beyond removing a unit and
annotating a label, and the question of alias tolerance that our tooling raised but
did not test. A further limitation concerns the units experiment specifically. Unlike the
label collision, whose prevalence we measured exhaustively, we did not measure how
often a real interface or export drops the unit column, so we treat the units experiment
as a controlled demonstration of the genuinely absent case rather than as evidence
about its frequency in the wild. We invite the community to run the task sets on their own agents and
report where, for their clients, the extra cost stops being worth paying.

\paragraph{Threats to validity.}
A broken change in this setting returns fluent, plausible prose, so several of our
own early conclusions were wrong until we read the bytes. We record the main cases,
because they also motivate the checks in Section 3. First, a single target units
pilot suggested the lower-performing agent could not convert $\mu$M to nM at all; the 20
target run showed it converts on 28 percent of rows, so the pilot floor was a
property of one table. Second, two defects in the numeric scorer (a micro sign
variant it did not recognize, and a rule that captured a secondary measurement)
inflated an apparent control failure rate about tenfold; both were fixed and the
few genuinely ambiguous answers were adjudicated by hand. Third, a response window
that was not reproducible, together with an underpowered 30-item label subset,
first suggested that the label annotation had lost its effect. A reproducible run
over the full 200-row window restored a strong effect ($p \approx 4\times10^{-16}$).
The same run also corrected an earlier claim that the annotation closed the gap
between the two agents. It narrows the gap but does not close it. These
episodes are the reason the study fixes windows by ordering, checks payload
identity per item, and verifies every number against the raw bytes.

\section{Conclusion}
 
We asked whether making a tool's return value more explicit helps an agent, and at
what token cost, by changing one column of what a production life-science server
returns while holding the model and the task fixed. The answer is that the value of
explicitness is a property of the client, not of the payload. When the needed
information was present but implicit, a zero-information column lifted a
lower-performing agent from 16 to 75 of 124 correct, for about 407 tokens, under one
percent of the run. A higher-performing agent already answered 123 of 124 and gained
almost nothing. The column narrowed the gap between the two agents but did not close
it. When the information was genuinely absent, removing the unit made the
higher-performing agent wrong on 350 of 929 micromolar rows and the lower-performing
agent wrong on all of them, and it raised token cost several times even on the rows
that stayed correct. Because these collisions are the common case rather than the
exception, and because a public server cannot see which agent is calling, a uniform
instruction to be explicit is best read as a safeguard: cheap, worthless to the
clients that would have inferred the answer, and decisive for those that would not.
The framework we release lets an operator measure, for a given set of clients,
where that extra cost stops being worth paying, and the question it makes measurable,
namely the efficiency of knowledge sharing between a language model and an external
source, will outlast the particular protocol studied here.

\section{Acknowledgments}
This work was supported by the MEXT National Life Science Database Project (NLDP) (grant number JPNLDP202401) and the Life Science Database Integration Project, NBDC of Japan Science and Technology Agency.
We also acknowledge and thank the participants of the DBCLS Domestic BioHackathon BH26.6, with whom several of the authors had fruitful discussions and gained deeper insights into this work.

\section*{Declaration on Generative AI}
  
 During the preparation of this work, the authors used Claude Desktop in order to: Drafting content. After using this tool, the authors reviewed and edited the content as needed and take full responsibility for the publication’s content. 
 
\bibliography{references}

\section{Online Resources}

The sources to reproduce this work and result data are available via
\href{https://github.com/dbcls/mcp-return-value-study}{GitHub}.

\end{document}
